% Lean: EconLean/Preferences/{Basic,Utility}.lean
\chapter{Preferences}
\label{chap:preferences}

\section{Preference Relations and Rationality}
\label{sec:preferences}

This section corresponds to \texttt{EconLean/Preferences/Basic.lean} and formalizes
Mas-Colell, Whinston \& Green (MWG) \S1.B (Proposition~1.B.1). Let $X$ be a set of
alternatives and $R \subseteq X \times X$ a binary relation written $x \pref y$
(``$x$ is at least as good as $y$'').

\begin{definition}[Strict preference]
\label{def:P}
\lean{Econ.P}
\leanok
Given a weak preference relation $R$ on $X$, the \emph{strict preference} relation
$P(R)$ (written $x \spref y$) is defined by
\[
  x \spref y \iff x \pref y \land \neg(y \pref x).
\]
\end{definition}

\begin{definition}[Indifference]
\label{def:I}
\lean{Econ.I}
\leanok
Given a weak preference relation $R$ on $X$, the \emph{indifference} relation
$I(R)$ (written $x \indiff y$) is defined by
\[
  x \indiff y \iff x \pref y \land y \pref x.
\]
\end{definition}

\begin{definition}[Rational preference]
\label{def:rational}
\lean{Econ.Rational}
\leanok
A preference relation $R$ on $X$ is \emph{rational} if it is:
\begin{enumerate}
  \item \emph{complete}: for all $x, y \in X$, $x \pref y$ or $y \pref x$;
  \item \emph{transitive}: for all $x, y, z \in X$, if $x \pref y$ and $y \pref z$, then $x \pref z$.
\end{enumerate}
\end{definition}

\begin{proposition}[Reflexivity]
\label{thm:refl}
\lean{Econ.Rational.refl}
\leanok
\uses{def:rational}
If $R$ is rational, then $R$ is reflexive: $x \pref x$ for every $x \in X$.
\end{proposition}
\begin{proof}
\leanok
By completeness applied to $x$ and $x$, either $x \pref x$ or $x \pref x$; in both cases $x \pref x$.
\end{proof}

\begin{lemma}[Strict preference via negation]
\label{lem:P_iff_not}
\lean{Econ.Rational.P_iff_not}
\leanok
\uses{def:P, def:rational}
If $R$ is rational, then for all $x, y \in X$,
\[
  x \spref y \iff \neg(y \pref x).
\]
\end{lemma}
\begin{proof}
\leanok
The forward direction is immediate from Definition~\ref{def:P}. Conversely, if
$\neg(y \pref x)$, completeness yields $x \pref y \lor y \pref x$, so $x \pref y$
must hold, giving $x \spref y$.
\end{proof}

\begin{proposition}[Properties of strict preference]
\label{thm:P_props}
\lean{Econ.P_irrefl, Econ.P_trans, Econ.P_of_P_of_R, Econ.P_of_R_of_P}
\leanok
\uses{def:P}
Let $R$ be a binary relation on $X$.
\begin{enumerate}
  \item \emph{Irreflexivity}: for all $x \in X$, $\neg(x \spref x)$ (no hypotheses on $R$ needed).
  \item If $R$ is transitive, then $\spref$ is transitive: $x \spref y$ and $y \spref z$ imply $x \spref z$.
  \item \emph{Mixed transitivity}: if $R$ is transitive, then
        $(x \spref y \land y \pref z) \implies x \spref z$ and
        $(x \pref y \land y \spref z) \implies x \spref z$.
\end{enumerate}
\end{proposition}
\begin{proof}
\leanok
For (1), $x \spref x$ means $x \pref x \land \neg(x \pref x)$, a contradiction.
For (2), $x \pref y \pref z$ gives $x \pref z$ by transitivity; if $z \pref x$, then
$z \pref x \pref y$ would imply $z \pref y$, contradicting $y \spref z$.
For (3), if $x \spref y \pref z$, transitivity gives $x \pref z$, and $z \pref x$
would yield $y \pref z \pref x$, contradicting $x \spref y$. The symmetric case
$x \pref y \spref z$ is analogous.
\end{proof}

\begin{proposition}[Indifference is an equivalence relation]
\label{thm:I_equiv}
\lean{Econ.Rational.I_equivalence}
\leanok
\uses{def:I, def:rational}
If $R$ is rational, then $\indiff$ is an equivalence relation on $X$.
\end{proposition}
\begin{proof}
\uses{thm:refl}
\leanok
Reflexivity follows from $x \pref x$ (Proposition~\ref{thm:refl}). Symmetry is immediate
from the symmetry of $x \pref y \land y \pref x$. For transitivity, $x \indiff y$ and
$y \indiff z$ yield $x \pref y \pref z$ and $z \pref y \pref x$, so $x \indiff z$.
\end{proof}

\begin{proposition}[Preorder structure]
\label{thm:isPreorder}
\lean{Econ.Rational.isPreorder}
\leanok
\uses{def:rational}
Every rational preference relation $R$ is a preorder (reflexive and transitive).
\end{proposition}
\begin{proof}
\uses{thm:refl}
\leanok
Reflexivity is Proposition~\ref{thm:refl}, and transitivity is part of Definition~\ref{def:rational}.
\end{proof}

\begin{proposition}[Trichotomy]
\label{thm:trichotomy}
\lean{Econ.Rational.trichotomy}
\leanok
\uses{def:P, def:I, def:rational}
If $R$ is rational, then for all $x, y \in X$, at least one of $x \spref y$,
$x \indiff y$, or $y \spref x$ holds (and by definition they are pairwise exclusive).
\end{proposition}
\begin{proof}
\leanok
Case-split on whether $x \pref y$ and $y \pref x$ hold. The three combinations where at
least one holds give $x \indiff y$, $x \spref y$, or $y \spref x$; the fourth case
($\neg(x \pref y)$ and $\neg(y \pref x)$) contradicts completeness.
\end{proof}


\section{Utility Representations}
\label{sec:utility}

This section corresponds to \texttt{EconLean/Preferences/Utility.lean} and formalizes
MWG \S1.B (Definition~1.B.2, Proposition~1.B.2) and the finite case of \S3.C.

\begin{definition}[Utility representation]
\label{def:represents}
\lean{Econ.Represents}
\leanok
A function $u : X \to \R$ \emph{represents} a preference relation $R$ on $X$ if
\[
  \forall x, y \in X,\quad x \pref y \iff u(y) \le u(x).
\]
\end{definition}

\begin{theorem}[Representable preferences are rational; MWG Prop.~1.B.2]
\label{thm:represents_rational}
\lean{Econ.Represents.rational}
\leanok
\uses{def:represents, def:rational}
If $u : X \to \R$ represents $R$, then $R$ is rational.
\end{theorem}
\begin{proof}
\leanok
Completeness: for any $x, y \in X$, totality of $\le$ on $\R$ gives $u(y) \le u(x)$ or
$u(x) \le u(y)$, hence $x \pref y$ or $y \pref x$. Transitivity: if $x \pref y$ and
$y \pref z$, then $u(z) \le u(y) \le u(x)$, so $u(z) \le u(x)$ and thus $x \pref z$.
\end{proof}

\begin{lemma}[Strict preference under utility representation]
\label{lem:represents_P_iff}
\lean{Econ.Represents.P_iff}
\leanok
\uses{def:represents, def:P}
If $u$ represents $R$, then for all $x, y \in X$,
\[
  x \spref y \iff u(y) < u(x).
\]
\end{lemma}
\begin{proof}
\leanok
Unfolding $P$, $x \spref y \iff u(y) \le u(x) \land \neg(u(x) \le u(y)) \iff u(y) < u(x)$.
\end{proof}

\begin{proposition}[Ordinal invariance of utility]
\label{thm:comp_strictMono}
\lean{Econ.Represents.comp_strictMono}
\leanok
\uses{def:represents}
If $u : X \to \R$ represents $R$ and $f : \R \to \R$ is strictly increasing, then
$f \circ u$ also represents $R$.
\end{proposition}
\begin{proof}
\leanok
Because $f$ is strictly increasing, $f(u(y)) \le f(u(x)) \iff u(y) \le u(x) \iff x \pref y$.
\end{proof}

\begin{definition}[Preference induced by a utility function]
\label{def:prefOf}
\lean{Econ.prefOf, Econ.represents_prefOf}
\leanok
\uses{def:represents}
Any function $u : X \to \R$ induces a preference relation $\pref_u$ defined by
$x \pref_u y \iff u(y) \le u(x)$, which is represented by $u$ tautologically.
\end{definition}

\begin{definition}[Lower contour set and counting utility]
\label{def:countingUtility}
\lean{Econ.lowerContour, Econ.mem_lowerContour, Econ.countingUtility}
\leanok
Suppose $X$ is finite. For each $x \in X$, define its \emph{lower contour set}
\[
  L(x) = \{y \in X \mid x \pref y\}
\]
and the \emph{counting utility} $u(x) = |L(x)| \in \R$.
\end{definition}

\begin{lemma}[Monotonicity of lower contour sets]
\label{lem:lowerContour_subset}
\lean{Econ.lowerContour_subset}
\leanok
\uses{def:rational, def:countingUtility}
If $X$ is finite, $R$ is rational, and $x \pref y$, then $L(y) \subseteq L(x)$.
\end{lemma}
\begin{proof}
\leanok
If $z \in L(y)$, then $y \pref z$; by transitivity with $x \pref y$, we obtain
$x \pref z$, so $z \in L(x)$.
\end{proof}

\begin{theorem}[Existence of utility on finite sets]
\label{thm:exists_utility_of_finite}
\lean{Econ.exists_utility_of_finite}
\leanok
\uses{def:rational, def:represents}
If $X$ is finite and $R$ is rational, then there exists $u : X \to \R$ representing $R$.
\end{theorem}
\begin{proof}
\uses{def:countingUtility, lem:lowerContour_subset, thm:refl}
\leanok
Take $u(x) = |L(x)|$, the counting utility.
\begin{itemize}
  \item[$(\Rightarrow)$] If $x \pref y$, Lemma~\ref{lem:lowerContour_subset} gives
        $L(y) \subseteq L(x)$, hence $|L(y)| \le |L(x)|$, i.e.\ $u(y) \le u(x)$.
  \item[$(\Leftarrow)$] Suppose $u(y) \le u(x)$ but $\neg(x \pref y)$. By completeness,
        $y \pref x$, so $L(x) \subseteq L(y)$ by Lemma~\ref{lem:lowerContour_subset}.
        Moreover $y \in L(y)$ by reflexivity (Proposition~\ref{thm:refl}) while
        $y \notin L(x)$, so $L(x) \subsetneq L(y)$. Hence $|L(x)| < |L(y)|$,
        contradicting $u(y) \le u(x)$.
\end{itemize}
\end{proof}

\begin{corollary}[Characterization of rationality on finite sets]
\label{thm:rational_iff_utility}
\lean{Econ.rational_iff_exists_utility}
\leanok
\uses{def:rational, def:represents}
On a finite set $X$, a preference relation $R$ is rational if and only if it admits a
utility representation.
\end{corollary}
\begin{proof}
\uses{thm:exists_utility_of_finite, thm:represents_rational}
\leanok
Combine Theorem~\ref{thm:exists_utility_of_finite} and Theorem~\ref{thm:represents_rational}.
\end{proof}


